\section{Methods}\label{sec:methods}

\subsection{The \NMethodsWord{} methods}
MICA is the method under test. Two baselines are the ones that its paper compares with, ModelAngelo and EModelX(+AF); four are newer deep-learning model builders, EMProt, E3-CryoFold, CryoAtom and CryoZeta. EMProt and E3-CryoFold passed a published screen (Section~\ref{sec:screen}); CryoAtom, which had passed it, and CryoZeta, which had been kept in reserve, were added after the first version of this report (Section~\ref{sec:added}). Every method ran with the settings of the example in its own README, with no per-map tuning. Table~\ref{tab:methods} gives the version, the inputs and the file that was scored for each. All methods that accept predicted structures got the same AlphaFold~3 (AF3) predictions; ModelAngelo, E3-CryoFold, CryoAtom and CryoZeta, which have no such input, got the map and the sequence only (CryoZeta also an alignment of every chain).

\begin{table}[t]
\centering\small
\caption{The methods and how they were run. All commits are pinned in \texttt{methods/methods\_table.csv} and in the run information of each session. The scored file is the one fixed before any score was seen (Section~\ref{sec:metrics}). $^\dagger$\,Added after the first version of this report (Section~\ref{sec:added}).}
\label{tab:methods}
\begin{tabularx}{\textwidth}{@{}>{\raggedright\arraybackslash}p{2.45cm}>{\raggedright\arraybackslash}p{1.95cm}>{\raggedright\arraybackslash}p{1.85cm}>{\raggedright\arraybackslash}X@{}}
\toprule
Method & Paper & Commit & Inputs used here; file that is scored \\
\midrule
MICA & \citet{gyawali2025mica} & \texttt{\CommitMica} & Map and one AF3 prediction per unique chain, docked into the map by MICA's program 2 (Section~\ref{sec:mica}). Scored: \texttt{MICA\_all\_atom\_model.pdb}. \\
ModelAngelo & \citet{jamali2024modelangelo} & \texttt{\CommitAngelo} & Map and sequence; no AlphaFold input (\texttt{model\_angelo build}). Scored: the sequence-assigned model. \\
EModelX(+AF) & \citet{chen2024emodelx} & \texttt{\CommitEmodelx} & Map, sequence and the AF3 predictions as templates (protocol \texttt{temp\_flex}, no AlphaFold database download). Scored: \texttt{all\_atom\_model.pdb} (after Pulchra). \\
EMProt & \citet{li2026emprot} & \texttt{\CommitEmprot} & Map, sequence and the AF3 predictions as per-chain models (\texttt{emprot build ... --chain}). Scored: \texttt{output.cif}. \\
E3-CryoFold & \citet{wang2025e3cryofold} & \texttt{\CommitEthree} & Map and sequence (\texttt{inference.py --protocol pre\_align --t 0.1}); no AlphaFold input. Scored: the \texttt{pre\_align} model. \\
CryoAtom$^\dagger$ & \citet{su2026cryoatom} & \texttt{\CommitCryoatom} & Map and sequence (\texttt{python -m CryoAtom build -s FASTA -v MAP -o OUT}, v1.0.0); no AlphaFold input. Scored: the model \texttt{<name>.cif} (after pruning). \\
CryoZeta$^\dagger$ & \citet{zhang2026cryozeta} & \texttt{\CommitCryozeta} & Map, sequence and a multiple sequence alignment (MSA) of every unique chain from the ColabFold server (\texttt{cryozeta-detection}, then \texttt{cryozeta-inference}, settings of the demo); no AlphaFold input. Scored: sample~0, the best of five samples (\texttt{<name>\_sample\_0.cif}). \\
\bottomrule
\end{tabularx}
\end{table}

\subsection{MICA and how it was run here}\label{sec:mica}
MICA combines the cryo-EM map with AF3 predictions in three programs \citep{gyawali2025mica}. Program~1 splits every AF3 prediction into domains (Merizo; \citealp{lau2023merizo}). Program~2 docks each domain into the map with Phenix \texttt{dock\_in\_map} \citep{liebschner2019phenix}. Program~3 resamples the map to 1~\AA{} voxels, predicts maps window by window with a 3D neural network (windows of $64^3$ voxels, step 48), which takes the docked AF3 models as extra input channels, builds the chains from these predictions and from the AF3 structures, completes the side chains with Pulchra \citep{rotkiewicz2008pulchra} and refines the model against the map with Phenix real-space refinement \citep{afonine2018rsr}. MICA's code and weights are used as released (commit \texttt{\CommitMica}); the network, the building of the chains and the refinement are not changed. Four things differ from the way the authors run MICA, each forced by free compute; the evidence is below or in Section~\ref{sec:dockprocedure}, and the effects are discussed in Sections~\ref{sec:dockresults} and~\ref{sec:limits}:
\begin{enumerate}
\item \textbf{Docking.} On a test entry (26DK, one chain of 191 residues, 3.5~\AA{}) MICA's own docking call (Phenix molecular replacement, no time limit) did not finish in 66 minutes on a 4-core, 31~GB Kaggle machine and was then killed by the system for lack of RAM. Program~2 is therefore run through a small driver that imports MICA's own file and replaces only the method that calls Phenix: \texttt{dock\_in\_map} in its Fourier-search mode with \texttt{quick=True}, at 6~\AA{} (then 10~\AA{}), with a time limit per domain. The placements are coarse (Section~\ref{sec:dockresults}).
\item \textbf{No docked input where nothing can be docked.} If no domain of an entry docks, MICA still runs without docked AF3 input, which its README describes as an option (``skip this step if you don't want to use AF3 at input level''); the AF3 channels of the network are then zero. The AF3 structures are still used when MICA builds the chains. Every MICA result is labelled by its docking mode: \emph{full} (every domain docked), \emph{partial} or \emph{none}.
\item \textbf{Program 3 on a GPU, one window at a time.} MICA predicts windows in batches of up to 8 when a map has more than 200 windows. That setting (\texttt{batch\_threshold = 200} in \texttt{utils/predict.py}) exceeded the 15~GB of a T4 on the first three large maps. In the GPU sessions of the second and third round it is raised to $10^6$, so that every window is predicted alone, as MICA does for small maps. This changes memory use only.
\item \textbf{No real-space refinement.} MICA's last step, Phenix real-space refinement with the parameter file shipped in MICA's repository (written for Phenix 1.20.1), stops with a parameter error in the only Phenix version available here (\PhenixVersion: ``Incompatible parameter objects: scope \texttt{cdl\_nucleotides}''). It stopped with this error in every one of the 25 program-3 runs of the GPU sessions that built a model (the main sessions and the session without docked input; the run logs are in the repository), while the program itself reports exit code~0. The scored file is therefore MICA's all-atom model after Pulchra and before refinement, chosen before any MICA result existed because it is the program's own output. Refinement moves atoms only a little and is not expected to change the TM-score much, but its effect on the other metrics is not measured (Section~\ref{sec:limits}).
\end{enumerate}

\subsection{The baselines}
\emph{ModelAngelo} \citep{jamali2024modelangelo} was run as \texttt{model\_angelo build -v MAP -pf FASTA -o OUT --device 0} with the default weights bundle. \emph{EModelX(+AF)} \citep{chen2024emodelx} was run in its AlphaFold mode (\texttt{--protocol=temp\_flex}) with the AF3 predictions as the templates of its README example (one folder per sequence with \texttt{ranked\_0.pdb}), followed by Pulchra. Feeding AF3 models to it is a protocol choice made so that all methods with an AlphaFold input get the same predictions. \emph{EMProt} \citep{li2026emprot} was run as \texttt{emprot build --map MAP --seq FASTA --chain <AF3 models> --device 0}. \emph{E3-CryoFold} \citep{wang2025e3cryofold} was run with the \texttt{pre\_align} protocol of its README example; its \texttt{denovo} protocol is not part of the benchmark. \emph{CryoAtom} \citep{su2026cryoatom} (v1.0.0, the code of the paper; its successor CryoAtom2 is a preprint and was not used) was run as \texttt{python -m CryoAtom build -s FASTA -v MAP -o OUT}, the command of its README, with the weights of the authors' server and the ESM-2 language model that the program loads; it has no AlphaFold input. \emph{CryoZeta} \citep{zhang2026cryozeta} (no release tag; its weights are free for academic and non-commercial research only) was run with the two programs of its README: \texttt{cryozeta-detection}, a 3D network that finds C$\alpha$ positions in the map (the contour level is the one recommended in the EMDB header), and \texttt{cryozeta-inference}, a diffusion model that builds the structure from the sequence, an MSA of every unique chain and these detections. The MSAs came from the public ColabFold server (the README names the ColabFold search as the way to get them); the server's alignments carry no species, so the chains of a complex are not paired across sequences. The settings are those of the demo of the authors (seed 101, 5~samples, 20 diffusion steps, 10 EMPairformer cycles, bfloat16 numbers). The compiled kernels of the demo script (DeepSpeed and cuEquivariance attention, a fast layer norm) need a CUDA compiler and an Ampere GPU and are off; the program has a plain implementation of each. The program ranks its five samples and writes the best as sample~0, which is the model of the benchmark; the other four are scored next to it (Appendix, Table~\ref{tab:variants}). No line of its source was changed.

\subsection{Selection of the newer methods}\label{sec:screen}
The newer methods were chosen with a screen that was written down before the search (spec decision~2): a method had to be published after both baselines, use deep learning, build protein atomic models from a map and a sequence, have public code and weights, run on a free 16~GB GPU and have documented defaults. A literature and repository search (to 11~October 2026) screened \NScreened{} methods: \NScreenPassed{} passed (CryoAtom, EMProt and E3-CryoFold), \NScreenReserve{} were kept in reserve (EM3DFold, a preprint; CryoZeta, which needs 32~GB of GPU memory) and \NScreenExcluded{} were excluded, each with a reason (Appendix, Table~\ref{tab:screen}). CryoAtom \citep{su2026cryoatom} was first dropped before the production runs, as the first scope cut of the spec allows: in a pilot on the smallest entry neither of its two versions finished within 2 hours, and at the floor of 12 entries it would have needed more GPU hours than one weekly quota. The first version of this report therefore had two newer methods, not three. The decision did not depend on any result. CryoAtom and CryoZeta were added afterwards at the request of Prof.\ Cheng (Section~\ref{sec:added}). Forward citations of MICA could not be searched (the citation databases were not reachable from the workspace), so the screen may have missed methods that only a citation search finds.

\subsection{Methods added after the first version}\label{sec:added}
After the first version of this report, Prof.\ Cheng asked me (31~October 2026) to add CryoAtom and CryoZeta to the benchmark. I added them and did not replace any method, and I wrote the decisions into the project log before any run of either method. The twenty entries, the AF3 inputs, the scoring protocol, the measures and the statistical tests are unchanged; the one change of the analysis is that the Holm correction covers the pairs of all added methods with MICA (\NFamilyWord{} pairs in this version, four in the first). The first version of the report stays as it was sent. Four rules are specific to the added methods:
\begin{itemize}
\item \textbf{Run limit.} 10~hours per run for CryoAtom, against 3~hours for the other baselines, because its pilot did not finish within 2~hours, and 6~hours for CryoZeta, half of a 12-hour session. The longer limits are a favour to both methods: \NOverThreeHCryoatom{} of the CryoAtom models (\OverThreeHListCryoatom{}) and \NOverThreeHCryozeta{} of the CryoZeta models (\OverThreeHListCryozeta{}) took longer than the 3~hours of the other baselines and would have been failures under their limit. No CryoAtom run reached 10~hours; the limit of 6~hours stopped \FailTimeoutCryozeta{} of the CryoZeta runs (\FailTimeoutListCryozeta{}).
\item \textbf{GPU memory.} A CUDA out-of-memory error on the 15~GB GPU is a failed run of the method on this hardware and not a fault of the installation (the README of CryoAtom asks for 13 to 14~GB, that of CryoZeta for 32~GB); the notebook records it as such, and the entries stay in every table. No larger GPU was used, so these entries were not run again. CryoZeta starts with its default attention chunk size of 256 and, after an out-of-memory error, tries 64, 16 and 4, a setting that changes the memory and the time but not the algorithm; the entry is failed only when all four attempts have failed. A run that stops for another reason is repeated once in float32.
\item \textbf{One line of the source of CryoAtom is changed}, so that two programs can run at the same time (one per GPU, as for the other methods): the program writes the profile of every chain to the fixed file \texttt{/tmp/CryoAtom\_temp.hmm}, which two programs would overwrite and delete, and the change adds the process number to the file name. The notebook checks that the line occurs exactly once, and the change is saved as \texttt{cryoatom\_source\_patch.diff} in the run record. Nothing else of the program is changed.
\item \textbf{CryoZeta ran on as many entries as the free GPU time allowed.} A CryoZeta run takes from \RtMinCryozeta{}~min to several hours on the T4 (Table~\ref{tab:runtime}), and the free machine allowed four sessions with two programs at a time, in which not all twenty entries fit. The entries ran in an order that was written down before they had a result: in the order of the entry list in the first session and, for the entries still without a result, the smallest deposited model first in the second and the third (an entry that the system had killed went first and alone, and the entry that an earlier session had cut last). The list of the fourth session was set by hand to 9VUH, 9WFL and 9WME, with 9VUH first and alone; 9VI4 was not in it. A run that the system kills while the other GPU is busy is run again alone, as for the other baselines, and a kill while alone is a failure (26DK, in the detection step of the program, in the first two sessions, and 9OC0 in the same step in the third). An entry that was not started, was cut by the end of a session or was killed without time for a second run is \emph{not run}: it is not a failure of the method, it is left out of the comparisons and it is marked N in the tables. The results of CryoZeta are therefore those of the \NModelCryozeta{} entries that it finished, a subset that is not random; it was not run on \NotRunListCryozeta{}.
\end{itemize}
CryoAtom ran in a notebook of its own (the scheduler, the two GPUs and the scoring are those of the other methods) on all \NEntries{} entries in one GPU session, and was scored in a separate CPU session with the same code. CryoZeta ran in the same kind of notebook in four GPU sessions that ran entries; an earlier one ended with a set-up error before any entry was run (two data files of the program had not been downloaded; the notebook was corrected), and the first start of the third stopped after 25~seconds at the check of the internet (one TLS time-out; the check now asks again). Its models were scored in two CPU sessions; the two scoring sessions that followed the last two GPU sessions stopped by themselves, because those sessions had built no model (Table~\ref{tab:sessions}).
